% Tony Judt: Origins of the European Union: Reading Summary
% Introduction to the European Union, Week 2
% Compile with: pdflatex judt-origins-of-eu-reading-summary.tex

\documentclass[11pt,a4paper]{article}
\usepackage[utf8]{inputenc}
\usepackage[T1]{fontenc}
\usepackage[margin=2.5cm]{geometry}
\usepackage{booktabs}
\usepackage{enumitem}
\setlist{nosep,leftmargin=*}
\usepackage{parskip}
\usepackage{microtype}
\usepackage{hyperref}
\usepackage{xcolor}
\definecolor{eublue}{RGB}{0,51,153}

\hyphenation{Spitzenkandidat Spitzenkandidaten}
\hypersetup{
  colorlinks=true,
  linkcolor=eublue,
  urlcolor=eublue,
  pdftitle={Judt: Origins of the EU},
  pdfauthor={Introduction to the EU, UCM}
}

\begin{document}

\title{Tony Judt: Origins of the European Union\\From \textit{Postwar} (2005, pp. 153--161)}
\author{Introduction to the European Union\\Universidad Complutense de Madrid\\Week 2}
\date{}
\maketitle
\vspace{-1em}
\hrule
\vspace{1em}

\section*{Rhetorical Framework (Ethos, Logos, Pathos)}

\begin{table}[h]
\centering
\begin{tabular}{@{}lll@{}}
\toprule
\textbf{Appeal} & \textbf{Meaning} & \textbf{Judt's Use} \\
\midrule
\textbf{Ethos} & Credibility, authority, trust & Historian's voice; scholarly evidence; balanced, critical perspective \\
\textbf{Logos} & Logic, structure, argument & Chronological narrative; institutional analysis; pragmatic reasoning \\
\textbf{Pathos} & Emotion, stakes, persuasion & Postwar trauma; desire for peace; ``weakness as foundation'' \\
\bottomrule
\end{tabular}
\end{table}

\textbf{Key insight:} Judt emphasizes \textit{logos} (pragmatic causes, institutional logic) while building \textit{ethos} through careful historical argument. \textit{Pathos} appears in the backdrop---war, reconciliation, fear of repetition---rather than in overt emotional language.

\section*{Bibliographic Information}

\begin{itemize}
\item \textbf{Author:} Tony Judt
\item \textbf{Work:} \textit{Postwar: A History of Europe Since 1945}
\item \textbf{Pages:} 153--161
\item \textbf{Year:} 2005
\item \textbf{Chapter:} Origins of the European Union
\end{itemize}

\section*{Summary \& Key Points}

\subsection*{I. Pre-WWII Context}

\begin{itemize}
\item The idea of European union predates WWII, with customs unions and pan-European visions emerging in the 19th and early 20th centuries.
\item WWI intensified calls for European unity, with leaders like Aristide Briand and Charles Gide advocating for overcoming rivalries and economic collaboration.
\item The interwar period saw attempts at economic cooperation (Steel Pact, customs unions), but these were limited by crises and national interests.
\end{itemize}

\textbf{Rhetorical note:} Judt establishes \textit{logos}---a causal chain (ideas $\to$ WWI $\to$ interwar attempts $\to$ failure) that sets up why postwar integration would differ.

\subsection*{II. Post-WWII Developments}

\begin{itemize}
\item WWII and Nazi occupation paradoxically unified much of Europe technically, but under authoritarian rule. Postwar, the desire for a peaceful, collaborative Europe grew stronger.
\item Churchill, Monnet, and Schuman were key advocates for European unity, emphasizing partnership between France and Germany as foundational.
\item The Council of Europe (1949) was an early forum for European conversation, but lacked real power.
\item The Benelux Agreement (1948) was a successful customs union among Belgium, Luxembourg, and the Netherlands.
\item The Marshall Plan and postwar recovery highlighted the need to integrate Germany into the European economy, balancing security and prosperity.
\end{itemize}

\textbf{Rhetorical note:}
\begin{itemize}
\item \textbf{Ethos:} Named figures (Churchill, Monnet, Schuman) lend authority.
\item \textbf{Pathos:} ``Desire for peaceful, collaborative Europe'' taps into postwar trauma and hope.
\item \textbf{Logos:} Institutional sequence (Council of Europe $\to$ Benelux $\to$ Marshall Plan) shows step-by-step logic.
\end{itemize}

\subsection*{III. The Schuman Plan and ECSC}

\begin{itemize}
\item The Schuman Plan (1950) proposed joint management of French and German coal and steel, leading to the ECSC (1951) with six founding states.
\item The ECSC was a political solution to Franco-German hostility, not just an economic project.
\item Christian Democratic leaders from border regions played a key role.
\item Britain declined to join the ECSC and later the EEC, preferring national control and global trade links over supranational European integration.
\end{itemize}

\textbf{Rhetorical note:}
\begin{itemize}
\item \textbf{Logos:} Economic instrument (coal/steel) as political tool---cause-and-effect argument.
\item \textbf{Ethos:} ``Political solution'' frames Judt as interpreter of motives, not just chronicler.
\item \textbf{Pathos:} Franco-German hostility and its resolution carry emotional weight beneath the institutional story.
\end{itemize}

\subsection*{IV. Treaty of Rome and EEC}

\begin{itemize}
\item The Treaty of Rome (1957) established the EEC and Euratom, aiming for a common market and economic integration among the ECSC six.
\item The EEC was grounded in weakness, not strength, but mutual self-interest and the lessons of war drove cooperation.
\item The creation of the EEC marked a shift toward economic integration as the practical path for European unity.
\end{itemize}

\textbf{Rhetorical note:}
\begin{itemize}
\item \textbf{Pathos:} ``Grounded in weakness''---appeals to humility, shared need, postwar vulnerability.
\item \textbf{Logos:} ``Practical path'' vs.\ political union---strategic choice with clear reasoning.
\item \textbf{Ethos:} Judt's critical distance signals scholarly judgment.
\end{itemize}

\section*{Main Arguments}

\begin{enumerate}
\item \textbf{Pragmatism over Idealism} --- European integration driven by practical concerns, not just visionary ideals.
\item \textbf{Franco-German Reconciliation} --- Partnership as moral and political foundation.
\item \textbf{Economic Integration as Path} --- Economic cooperation chosen as practical route to European unity.
\item \textbf{British Reluctance} --- Britain's different priorities; national control and global trade vs.\ supranational commitment.
\item \textbf{Weakness as Foundation} --- Mutual need, postwar vulnerability as driver of cooperation.
\end{enumerate}

\section*{Context \& Analysis}

\textbf{Judt's Approach:} Emphasizes the pragmatic, sometimes reluctant, steps toward European integration, shaped by security concerns, economic interests, and the need to overcome historical antagonisms.

\textbf{Rhetorical profile:} Heavy on \textit{logos} (pragmatism, institutions, causes); \textit{ethos} through scholarly tone and evidence; \textit{pathos} implicit in the stakes (peace, prosperity, avoiding another war).

\section*{Relevance to Course}

\begin{itemize}
\item \textbf{Ethos:} Historical context for why European integration began; Judt as authoritative source.
\item \textbf{Logos:} Pragmatic motivations, institutional choices, cause-and-effect.
\item \textbf{Pathos:} Stakes of peace, prosperity, and avoiding repetition of war.
\item \textbf{All three:} Critical perspective on founding---credibility, reasoning, and human stakes together.
\end{itemize}

\section*{Discussion Questions}

\begin{itemize}
\item How did WWII shape the desire for integration? (Ethos: whose experience? Logos: causal links? Pathos: trauma, fear, hope)
\item Why economic over political integration initially? (Ethos: who made the choice? Logos: constraints? Pathos: what was too risky?)
\item Pragmatism vs.\ idealism in founding? (Ethos: who framed it? Logos: evidence? Pathos: what did people want?)
\item Britain's decision not to join? (Ethos: British perspective; Logos: structural effects; Pathos: sovereignty, identity)
\item EEC ``grounded in weakness''? (Ethos: Judt's interpretation; Logos: how does weakness drive cooperation? Pathos: vulnerability, mutual need)
\end{itemize}

\vspace{1em}
\noindent\footnotesize\textit{Reference:} Judt, Tony. \textit{Postwar: A History of Europe Since 1945}. New York: Penguin Press, 2005, pp. 153--161.

\end{document}
